% !TeX TXS-program:compile = txs:///lualatex | txs:///view
\documentclass[12pt,a4paper]{article}
\usepackage[spanish]{babel}
\usepackage{amsmath}
\usepackage{amssymb}
\usepackage[top=2cm,bottom=2cm,left=2cm,right=2cm]{geometry}
\usepackage{amsthm}

\author{Carlos Alonso Aznarán Laos}
\title{Curso de \LaTeX \\ Práctica N$^{\circ}$2}
\newtheorem{ejer}{Ejercicio} % va en el preámbulo del documento (se llama declaración)

\begin{document}
\maketitle % todos los entornos empiezan con un begin y acaban en un end

\begin{ejer}
	Escribir matemática al interior de un texto.\\
	\noindent Por ejemplo, $\mathbb{N}=\left\{1,2,3,4,\dotsc\right\}$ es un
	conjunto infinito.\\

	Por ejemplo, \(\mathbb{N}=\left\{1,2,3,4,\dotsc\right\}\) es un conjunto
	infinito.\\

	Por ejemplo, \begin{math}\mathbb{N}=\left\{1,2,3,4,\dotsc\right\}\end{math}
	es un conjunto infinito.\\
\end{ejer}

\begin{ejer}
	Escribir
\end{ejer}

$H_{0}=\mu_{\text{con programa}}=11$.\\

$a^{m}\cdot a^{n}=a^{m+n}$.\\

${\left(\frac{a}{b}\right)}^{n}=\frac{{a}^{n}}{{b}^{n}}$, donde $b\neq0$.\\

$\sqrt[n]{ab}=\sqrt[n]{a}\cdot\sqrt[n]{b}$.\\

La propiedad $\sqrt[n]{\frac{a}{b}}=\frac{\sqrt[n]{a}}{\sqrt[n]{b}}$, donde $b\neq0$.\\
La propiedad $\displaystyle\sqrt[n]{\frac{a}{b}}=\frac{\sqrt[n]{a}}{\sqrt[n]{b}}$, donde $b\neq0$.\\

``Nuestro trabajo es, por lo tanto, mostrar que la integral\ldots
\begin{displaymath}
	\mathnormal{
		\int_{0}^{1}
		{\left(\widehat{f}\left(\alpha\right)\right)}^{3}
		\exp\left(\alpha n\right) d\alpha
	}
\end{displaymath} es distinta de cero.'' Helfgott, 2013: 714.
\end{document}
